% Глава 5: Результаты и анализ

\section{Результаты и анализ}
\label{sec:results}

\subsection{Диагностическая серия D (Phase~1)}
\label{sec:results_phase1}

Цель диагностической серии~--- проверить рабочую гипотезу: несут ли double-stream блоки
FLUX.1 стилевую информацию, достаточную для переноса стиля через LoRA-адаптер.
Серия включает эталонный запуск без адаптера (e00) и три конфигурации
с LoRA на double-stream блоках~[0–18]: базовую~(d01), с увеличенным числом шагов~(d02)
и с увеличенным рангом~(d03).

\subsubsection*{Количественные результаты}

\begin{table}[h]
\centering
\caption{Метрики диагностической серии D (стиль: Ван~Гог, 100~промптов MS-COCO)}
\label{tab:phase1_metrics}
\begin{tabular}{lcccccc}
\hline
\textbf{Эксперимент} & \textbf{Блоки} & \textbf{$r$} & \textbf{DINO-style\,$\uparrow$} & \textbf{CLIP-style\,$\uparrow$} & \textbf{FID\,$\downarrow$} & \textbf{LPIPS\,$\downarrow$} \\
\hline
e00 (без LoRA)           & —           & — & 0.046          & 0.463          & 283.3          & 0.840          \\
d01 (1\,000 шагов, $r{=}16$) & DS [0–18] & 16 & \textbf{0.191} & \textbf{0.486} & \textbf{246.3} & \textbf{0.790} \\
d02 (2\,000 шагов, $r{=}16$) & DS [0–18] & 16 & 0.162          & 0.483          & 257.5          & 0.790          \\
d03 (1\,000 шагов, $r{=}32$) & DS [0–18] & 32 & 0.102          & 0.467          & 271.8          & 0.810          \\
\hline
\end{tabular}
\end{table}

Конфигурация~d01 (ранг~$r{=}16$, 1\,000~шагов) демонстрирует наилучший результат
по всем ключевым метрикам. Метрика~DINO-style возросла с~$0.046$ (baseline) до~$0.191$,
то есть в~4{,}1~раза~--- что подтверждает гипотезу о том, что double-stream блоки
FLUX.1 несут стилевую информацию, доступную для извлечения через LoRA-дообучение.
Одновременно FID снизился с~283 до~246, а LPIPS~--- с~0.840 до~0.790, что свидетельствует
о повышении общего качества и стилевой близости генерируемых изображений.

Метрика~CLIP-content сохранилась на уровне~$0.254$--$0.255$ во всех конфигурациях
(включая baseline), что означает: применение LoRA не нарушает семантическую связь
сгенерированного изображения с текстовым промптом.

\subsubsection*{Влияние числа шагов и ранга}

Увеличение числа шагов с~1\,000 до~2\,000 (d02~vs~d01) приводит к снижению
DINO-style с~$0.191$ до~$0.162$ при сохранении LPIPS и ухудшении FID.
Этот эффект согласуется с известным явлением \textit{content overfitting} при длительном
дообучении LoRA в диффузионных моделях~\cite{frenkel2024blora}: адаптер начинает
«запоминать» структуру обучающего изображения в ущерб стилевому переносу.

Увеличение ранга с~16 до~32 при фиксированных 1\,000~шагах (d03~vs~d01) также
снижает DINO-style до~$0.102$. Предположительно, более высокий ранг допускает
большую ёмкость модели, что позволяет адаптеру захватывать как стилевые,
так и контентные характеристики обучающего изображения, тем самым размывая
стилевый сигнал в пространстве double-stream блоков.

\subsubsection*{Визуальная оценка}

При визуальном осмотре (10~изображений на конфигурацию) подтверждается
непоследовательность стилевого переноса в зависимости от семантики промпта.
Наружные сцены с большим числом объектов (толпа, пейзаж) воспроизводят
характерные для Ван~Гога мазки и тёплую охристую палитру.
Механические объекты (мотоцикл, спортивный инвентарь) сохраняют фотореалистичный
вид практически без следов стилизации.
Человеческая оценка качества d01: 3/5 (стиль виден на части изображений,
артефактов нет, семантика сохранена).

\subsubsection*{Предварительный вывод Phase~1}

По результатам диагностической серии~D подтверждена гипотеза о стилевой специализации
double-stream блоков FLUX.1; лучшая базовая конфигурация~--- $r{=}16$, 1\,000~шагов.
Окончательное решение GATE~1 принимается после серии~DC (Phase~1b, поиск Θ\_content).

% ----------------------------------------------------------------
\subsection{Серия DC (Phase~1b): поиск Θ\_content}
\label{sec:results_phase1b}

Цель серии~DC~--- проверить, можно ли выделить отдельный адаптер для контентных
признаков~(Θ\_content), тренируя LoRA исключительно на подмножестве double-stream блоков.
Три эксперимента варьируют диапазон блоков-кандидатов: ранние~[0–8], поздние~[9–18]
и наиболее значимые по итогам Phase~0 блоки~[0–5].

\subsubsection*{Количественные результаты}

\begin{table}[h]
\centering
\caption{Метрики серии DC~--- поиск Θ\_content (стиль: Ван~Гог, 100~промптов MS-COCO)}
\label{tab:phase1b_metrics}
\begin{tabular}{lcccc}
\hline
\textbf{Эксперимент} & \textbf{DINO-style\,$\uparrow$} & \textbf{CLIP-style\,$\uparrow$} & \textbf{FID\,$\downarrow$} & \textbf{LPIPS\,$\downarrow$} \\
\hline
dc\_content\_early\_ds (DS [0–8])  & 0.079          & 0.463          & 274.31          & 0.817          \\
dc\_content\_late\_ds  (DS [9–18]) & \textbf{0.121} & \textbf{0.480} & \textbf{257.22} & \textbf{0.817} \\
dc\_content\_from\_p0  (DS [0–5])  & 0.074          & 0.466          & 276.15          & 0.822          \\
\hline
\end{tabular}
\end{table}

Конфигурация~\texttt{dc\_content\_late\_ds} (блоки~[9–18]) превосходит остальных по всем
ключевым показателям: DINO-style достигает~$0.121$ против $0.079$ (ранние блоки)
и~$0.074$ (блоки~[0–5]), FID снижается до~$257$~--- наименьшее значение в серии.
Метрика CLIP-content у всех трёх конфигураций равна~${\approx}\,0.254$, что совпадает
с baseline без LoRA: стилевая адаптация не нарушает семантическую связь генерации
с текстовым промптом.

Конфигурация~\texttt{dc\_content\_from\_p0} (блоки~[0–5], без прогрева) показывает
наихудший FID и LPIPS, что свидетельствует о том, что обучение LoRA с нуля на
ограниченном числе ранних блоков снижает как общее качество изображений, так и
стилевую выразительность.

\subsubsection*{Вывод по GATE~1}

Серия~DC подтверждает, что поздние double-stream блоки~[9–18] несут наибольший
информационный сигнал, доступный для выделения в отдельный адаптер.
Оба условия GATE~1 выполнены: стилевые double-stream блоки~[0–18] подтверждены
в Phase~1, кандидат Θ\_content = блоки~[9–18] выделен в Phase~1b.

\textbf{Принято решение: Сценарий~A.}
Θ\_content~= double-stream~[9–18], Θ\_style~= double-stream~[0–18] (уточняется в Phase~2).
Это открывает Phase~6 (style-content mixing) и Phase~3.2 (раздельные~$\alpha_\mathrm{style}$ / $\alpha_\mathrm{content}$).

% ----------------------------------------------------------------
\subsection{Серия DA (Phase~2.1): диапазон double-stream блоков}
\label{sec:results_phase2_da}

Цель серии~DA~--- определить оптимальный диапазон double-stream блоков FLUX.1
для стилевого LoRA-адаптера. Фиксированы ранг~$r{=}32$ и 1\,000~шагов обучения;
варьируются граничные индексы блоков: $[0\text{–}6]$, $[0\text{–}12]$, $[0\text{–}18]$
и $[6\text{–}18]$. Конфигурация~$[0\text{–}18]$ (DA03) охватывает все 19~double-stream
блоков и служит верхней границей поиска.

\subsubsection*{Количественные результаты}

\begin{table}[h]
\centering
\caption{Метрики серии DA~--- диапазон DS-блоков ($r{=}32$, 1\,000~шагов, стиль: Ван~Гог)}
\label{tab:phase2_da_metrics}
\begin{tabular}{llccccc}
\hline
\textbf{Эксп.} & \textbf{Блоки} & \textbf{DINO-style\,$\uparrow$} & \textbf{CLIP-style\,$\uparrow$} & \textbf{CLIP-cont.\,$\uparrow$} & \textbf{FID\,$\downarrow$} & \textbf{LPIPS\,$\downarrow$} \\
\hline
DA01 & DS [0–6]  & 0.1007          & 0.4644          & 0.2503 & 270.55          & 0.8127          \\
DA02 & DS [0–12] & 0.1849          & 0.4973          & 0.2483 & 248.70          & 0.7879          \\
DA03 & DS [0–18] & \textbf{0.2568} & \textbf{0.5181} & 0.2447 & \textbf{235.85} & \textbf{0.7791} \\
DA04 & DS [6–18] & 0.1165          & 0.4774          & 0.2563 & 264.00          & 0.8170          \\
\hline
\end{tabular}
\end{table}

Конфигурация~DA03 (все double-stream блоки, $[0\text{–}18]$) лидирует по всем
пяти метрикам: DINO-style~$0.2568$~--- на 34\,\% выше лучшего результата
Phase~1 (d01: $0.191$); CLIP-style~$0.518$~--- наивысшее в серии; FID~$235.85$ и LPIPS~$0.779$~--- лучшие в серии.
Метрика CLIP-content стабильна во всех конфигурациях ($0.244$--$0.256$), что
подтверждает сохранение семантического контроля вне зависимости от диапазона блоков.

\subsubsection*{Вклад ранних и поздних блоков}

Сравнение DA04~($[6\text{–}18]$) с DA03~($[0\text{–}18]$) изолирует вклад блоков~$[0\text{–}5]$:
DINO-style падает с~$0.2568$ до~$0.1165$ (снижение в 2{,}2~раза),
FID ухудшается на~28 единиц.
Блоки~$[0\text{–}5]$ несут критически важный стилевой сигнал~--- это согласуется
с результатами Phase~0, где ds\_00 получил наивысший style specificity score ($+0.0175$).

Симметрично, конфигурация DA01~($[0\text{–}6]$) без поздних блоков показывает
DINO-style~$0.100$~--- столь же слабо, как и DA04. Стилизация возникает лишь при
достаточно широком покрытии, включающем как ранние, так и поздние блоки.
Добавление блоков~$[7\text{–}12]$ (DA01→DA02) даёт прирост DINO-style $+0.084$,
несмотря на то что по данным Phase~0 эти блоки content-dominant: по-видимому,
совокупный эффект широкого адаптера компенсирует низкую style specificity отдельных блоков.

\subsubsection*{Визуальная оценка и ограничение метода}

При визуальном осмотре (по 5~изображений на конфигурацию) DA03 получила наивысшую
среднюю человеческую оценку~$4{/}5$; DA01 и DA04 не показали видимой стилизации
(оценка $2{/}5$), DA02~--- частичную ($3{/}5$).

В ходе визуальной оценки обнаружена устойчивая зависимость качества стилизации
от семантической категории изображаемого объекта.
Современные предметы с сильными фотореалистичными prior-знаниями модели
(транспортные средства, готовая еда, спортивный инвентарь) практически не принимают
художественный стиль: стилевые признаки проявляются лишь как второстепенные
«отсылки»~--- например, в содержимом экрана или на фоновом знаке.
Напротив, объекты-контейнеры (экраны, вывески) и фоновые сцены принимают стиль
значительно свободнее, так как FLUX.1 не располагает жёстким фотографическим
prior-ом для их содержимого.

Этот феномен указывает на конкуренцию между текстовыми prior-ами модели
и стилевым сигналом LoRA-адаптера: при сильном prior-е FLUX.1 <<знает>>,
как должен выглядеть объект, и подавляет стилевую модификацию.
Количественно это выражается в относительно низких значениях DINO-style на промптах
с конкретными современными объектами при высоких значениях CLIP-content~---
семантика сохраняется, но стиль не переносится.
Данное ограничение характерно для метода в целом и подробнее рассмотрено
в разделе~\ref{sec:limitations}.

\subsubsection*{Вывод Phase~2.1}

Θ\_style~= double-stream~$[0\text{–}18]$ подтверждён как оптимальный диапазон.
Полный охват всех double-stream блоков даёт наилучшие количественные и визуальные
результаты; исключение любого подмножества (ранних или поздних блоков) ведёт к
значительному снижению DINO-style.
Данная конфигурация используется как фиксированная в следующих абляциях
(серии~DB: ранг~LoRA; DC: число шагов).
